\section[Conclusão]{Conclusão}

 Ao longo do desenvolvimento deste projeto, tive a oportunidade de crescer tanto tecnicamente quanto nas habilidades interpessoais. Acredito que evoluí muito tecnicamente, principalmente ao ter meu primeiro contato com desenvolvimento Front-End e fazer uso de diferentes linguagens e frameworks com os quais nunca tinha trabalhado antes. Também tive a oportunidade de familiarizar melhor com ferramentas de controle de versão, como o Git. Entretanto, percebo que ainda há muito espaço para aprimoramento técnico, uma vez que não consegui contribuir tanto com o desenvolvimento do código sem a ajuda de meus colegas. Nas habilidades de soft skills, observei uma grande melhora em minha comunicação e capacidade de trabalhar em equipe, devido a minha particiáção nas reuniões diárias, que contribuíram muito para meu desenvolvimento. Além disso, as situações adversas, como a enchente que impactou nossa equipe, reforçaram a importância da empatia e da flexibilidade. Foi crucial apoiar os colegas afetados e adaptar o plano de trabalho às circunstâncias imprevistas. Este projeto permitiu desenvolver um senso de resiliência e colaboração, características essenciais, muito valorizadas em qualquer ambiente profissional.
 
 O relacionamento entre as disciplinas cursadas e a AGES foi benéfico. As matérias anteriormente cursadas proporcionaram uma boa base em lógica de programação, que foi muito aplicada na prática durante o projeto no aprendizado de novas linguagens. Conceitos de engenharia de software, design de sistemas e gestão de projetos foram diretamente utilizados, mostrando a importância de um primeiro contato teórico com conceitos importantes antes de enfrentar desafios reais de desenvolvimento.
 
Este foi meu primeiro projeto na Agência e considero-o extremamente proveitoso, uma vez que esse apresentou uma série de desafios complexos que me forçaram a sair da zona de conforto e a encontrar soluções criativas. A diversidade das tarefas e a necessidade de um rápido desenvolvimento de habilidades técnicas e soft-skills, tornaram essa uma experiência abrangente e enriquecedora.

As lições aprendidas durante o projeto foram inúmeras. Destacam-se, entretanto, a capacidade de adaptação às mudanças, a necessidade de um planejamento detalhado, a necessidade de manter uma comunicação clara e constante com todos os membros da equipe e a importância de uma boa convivência e integração entre todos os integrantes do time. O envolvimento ativo com os \textit{stakeholders} também foi um ponto forte, garantindo que o projeto estivesse alinhado com as expectativas desde o início.

Durante todo o período de desenvolvimento do projeto, diversos pontos positivos foram evidenciados, destacando-se a determinação da equipe, que se mostrou empenhada em superar desafios; a escolha criteriosa das tecnologias utilizadas, como a linguagem Java para o backend e o framework Flutter para o \textit{frontend} (diferencial que facilitou uma boa implementação das funcionalidades previstas); e a utilização do Git como ferramenta de controle de versão, responsável por garantir um código limpo e organizado, permitindo um desenvolvimento mais eficiente e cooperativo. Além disso, a documentação bem estruturada e a criação de um protótipo detalhado no Figma contribuíram para a clareza, padronização e alinhamento do projeto com as expectativas dos \textit{stakeholders}. A capacidade da equipe em adaptar-se a contratempos, como a enchente que afetou o Rio Grande do Sul, demonstrou resiliência e compromisso com o projeto, além de empatia e solidariedade de todos os envolvidos. A proatividade e o apoio mútuo entre os membros da equipe foram fundamentais para a conclusão das tarefas e para a manutenção de um ambiente de trabalho positivo e produtivo.

Entretanto, Sempre há espaço para melhorar. Em termos de banco de dados, poderia ter havido uma melhor definição inicial da estrutura, evitando retrabalho. No desenvolvimento, uma maior atenção aos prazos e à distribuição equilibrada de tarefas teria prevenido sobrecargas e atrasos. 

Em retrospectiva, este projeto foi uma jornada de aprendizado intenso. As adversidades enfrentadas, desde contratempos técnicos, como minha falta de experiência inicial, até desastres naturais, reforçaram a importância de uma equipe colaborativa e resiliente. A experiência prática adquirida, aliada ao desenvolvimento pessoal, consolidou meu crescimento acadêmico e profissional. Estou confiante de que as lições aprendidas serão valiosas para futuros projetos, de modo que possa contribuir e evoluir cada vez mais. 
